\section{Conclusions}
\label{sec:conclusion}

We applied Topological Data Analysis to characterize the grokking phenomenon in a small transformer trained on modular arithmetic tasks. Across two datasets, five training fractions, and four architectural layers, we found that successful generalization is strongly associated with a structural regularization of the learned representation manifold. Specifically, $\beta_0$ converges to a stable value in both the decoder and linear layers concurrently with the rise in validation accuracy; i.e. decreasing in the decoder as intermediate representations consolidate into fewer connected components, and increasing in the linear layer as output representations organize into well-separated clusters; while intrinsic dimension decreases progressively across all layers to values substantially lower than those observed in non-generalizing conditions. These patterns are robust across all ten experimental conditions, and the partial generalization observed in the multiplication task at $q = 15\%$ produces intermediate topological behavior, suggesting the transition is continuous rather than abrupt. Future work should also test whether the identified signatures generalize across other cyclic groups $\mathbb{Z}_n$, non-commutative structures, and introducing topological regularization losses to accelerate or induce grokking by directly encouraging manifold simplification. 